\section{附录：工程落地检查清单}

\subsection{训练前准备清单}
为了把课程中的方法真正落地，团队通常需要先完成 ``训练前系统准备''。这一环节经常被低估，但它决定了后续 RL 是否能稳定推进。尤其是可验证 Agent 训练，任何日志缺失或工具观测缺失都会直接削弱反馈质量。

\begin{importantbox}{优先级排序建议}
先把 ``可观测'' 做完整，再追求 ``算法更强''。如果缺少高质量轨迹日志、工具调用回放和 verifier 诊断信息，再复杂的算法也会因为反馈噪声而收益有限。
\end{importantbox}

\begin{longtable}{p{2.6cm}p{3.4cm}p{5.0cm}}
\toprule
\textbf{阶段} & \textbf{检查项} & \textbf{通过标准} \\
\midrule
\endfirsthead
\toprule
\textbf{阶段} & \textbf{检查项} & \textbf{通过标准} \\
\midrule
\endhead
任务定义 & 任务目标可程序化表达 & 对每类任务可写出明确 success/failure 条件 \\
任务定义 & 终止条件明确 & 可判定 ``何时应停止继续调用工具'' \\
环境建模 & 环境状态可序列化 & 能回放每一步状态变更并定位差异 \\
环境建模 & 环境重置机制可靠 & 同一任务可重复运行并得到可比较结果 \\
工具链路 & 工具 I/O 协议固定 & 字段、类型、错误码、超时策略均有文档 \\
工具链路 & 工具副作用可审计 & 任意调用可追溯到操作者、参数、时间与结果 \\
Verifier & 结果层 verifier 完整 & 关键任务均有自动通过/失败判断 \\
Verifier & 过程层 verifier 可用 & 对关键中间步骤可检测常见违规模式 \\
数据管线 & 轨迹采集字段齐全 & Prompt、action、observation、reward 全量记录 \\
数据管线 & 失败样本可分型 & 至少区分工具失败、策略失败、验证器失败 \\
训练系统 & 采样吞吐可监控 & 可实时观察样本速率、成功率、拒绝率 \\
训练系统 & 更新稳定性可监控 & 可追踪梯度异常、loss 爆炸、熵突降等事件 \\
评估系统 & 难度分层报告可产出 & easy/medium/hard 分层结果可自动生成 \\
评估系统 & Harness 交换机制可运行 & 同任务可在不同 harness 下重复评测 \\
治理与安全 & 权限边界已实现 & 高风险工具调用需满足显式授权策略 \\
治理与安全 & 回滚链路可演练 & 关键失败可在分钟级回滚到安全状态 \\
\bottomrule
\end{longtable}

\subsection{训练中监控面板}
结合课程中的观察，团队在训练过程中至少应维护三块看板：能力看板（准确率/成功率）、探索看板（熵/轨迹去重率）、可靠性看板（格式正确率/工具错误率）。三者缺一不可。

\begin{knowledgebox}{推荐日常监控指标}
\begin{itemize}
    \item 任务成功率：按任务类别和难度分层统计，不看单一均值。
    \item 输出熵与去重率：监控探索是否过早塌缩。
    \item verifier disagreement：不同 verifier 对同轨迹是否冲突。
    \item 工具失败率：按工具类型和错误码分布聚类。
    \item 轨迹长度分布：防止出现无意义超长循环。
\end{itemize}
\end{knowledgebox}

\subsection{训练后验收与发布门槛}
课程强调 ``高分不等于可部署''。因此，训练后验收不能只看 benchmark 分数，还要看跨环境稳定性、异常场景恢复能力、权限边界遵守率。发布门槛必须写成制度化 checklist，而非临时判断。

\begin{warningbox}{发布前必须回答的三个问题}
\begin{itemize}
    \item 这个模型在分布外任务失败时，是否会安全失败（safe fail）？
    \item 当 verifier 冲突或不可用时，系统是否会降级而非盲目执行？
    \item 线上事故发生后，是否能在审计日志中还原完整责任链？
\end{itemize}
\end{warningbox}

\subsection{本章小结}
工程落地的关键不是 ``再加一种算法''，而是把任务定义、状态观测、验证体系、监控面板和发布门槛连接成闭环。只有闭环完整，课程中的方法才会在真实系统里持续生效。

